\documentclass[11pt,a4paper]{article}
\usepackage[utf8]{inputenc}
\usepackage[T1]{fontenc}
\usepackage{lmodern}
\usepackage{amsmath}
\usepackage[margin=2cm]{geometry}
\usepackage{enumitem}
\pagestyle{plain}
\setlength{\parindent}{0pt}
\setlength{\parskip}{0pt}
\begin{document}
{\Large\bfseries Tarea de hidruros}
\vspace{12pt}
\begin{enumerate}[label=\textbf{\arabic*.},leftmargin=*,itemsep=11pt,topsep=0pt]
\item \begin{minipage}[t]{\linewidth}
¿Cuál es la fórmula del hidruro de bario?

\vspace{5pt}
\begin{tabular}{@{}p{0.47\linewidth}p{0.47\linewidth}@{}}
a)~$\mathrm{BaH_{2}}$ & b)~$\mathrm{Ba_{2}H}$ \\
c)~BaH & d)~HBa \\
\end{tabular}
\end{minipage}

\item \begin{minipage}[t]{\linewidth}
¿Cuál es el nombre stock del $\mathrm{CoH_{3}}$?

\vspace{5pt}
\begin{tabular}{@{}p{0.47\linewidth}p{0.47\linewidth}@{}}
a)~hidruro de cobalto (II) & b)~trihidruro de cobalto (III) \\
c)~hidruro de cobalto (III) & d)~hidruro cobaltoso \\
\end{tabular}
\end{minipage}

\item \begin{minipage}[t]{\linewidth}
¿Cuál es el nombre tradicional del $\mathrm{NiH_{2}}$?

\vspace{5pt}
\begin{tabular}{@{}p{0.47\linewidth}p{0.47\linewidth}@{}}
a)~hidruro de níquel (III) & b)~hidruro niqueloso \\
c)~hidruro niquélico & d)~dihidruro niquélico \\
\end{tabular}
\end{minipage}

\item \begin{minipage}[t]{\linewidth}
¿Cuál es el nombre sistemático del $\mathrm{AlH_{3}}$?

\vspace{5pt}
\begin{tabular}{@{}p{0.47\linewidth}p{0.47\linewidth}@{}}
a)~trihidruro de aluminio & b)~monohidruro de aluminio \\
c)~trihidruro alumínico & d)~hidruro de aluminio (III) \\
\end{tabular}
\end{minipage}

\item \begin{minipage}[t]{\linewidth}
¿Cuál es la fórmula del hidruro cuproso?

\vspace{5pt}
\begin{tabular}{@{}p{0.47\linewidth}p{0.47\linewidth}@{}}
a)~CuH & b)~$\mathrm{CuH_{2}}$ \\
c)~$\mathrm{CuH_{3}}$ & d)~$\mathrm{Cu_{2}H}$ \\
\end{tabular}
\end{minipage}

\item \begin{minipage}[t]{\linewidth}
¿Cuál de estos compuestos es un hidruro volátil?

\vspace{5pt}
\begin{tabular}{@{}p{0.47\linewidth}p{0.47\linewidth}@{}}
a)~$\mathrm{MgH_{2}}$ & b)~$\mathrm{AlH_{3}}$ \\
c)~NaH & d)~$\mathrm{SiH_{4}}$ \\
\end{tabular}
\end{minipage}

\item \begin{minipage}[t]{\linewidth}
¿Cuál es la fórmula del metano?

\vspace{5pt}
\begin{tabular}{@{}p{0.47\linewidth}p{0.47\linewidth}@{}}
a)~$\mathrm{CH_{3}}$ & b)~$\mathrm{CH_{2}}$ \\
c)~$\mathrm{CH_{4}}$ & d)~$\mathrm{C_{4}H}$ \\
\end{tabular}
\end{minipage}

\item \begin{minipage}[t]{\linewidth}
¿Cuál es el nombre sistemático del $\mathrm{BH_{3}}$?

\vspace{5pt}
\begin{tabular}{@{}p{0.47\linewidth}p{0.47\linewidth}@{}}
a)~hidruro de boro (III) & b)~borano \\
c)~tetrahidruro de boro & d)~trihidruro de boro \\
\end{tabular}
\end{minipage}

\item \begin{minipage}[t]{\linewidth}
¿Cuál es la fórmula del hidruro plumboso?

\vspace{5pt}
\begin{tabular}{@{}p{0.47\linewidth}p{0.47\linewidth}@{}}
a)~$\mathrm{Pb_{2}H}$ & b)~$\mathrm{PbH_{2}}$ \\
c)~PbH & d)~$\mathrm{PbH_{4}}$ \\
\end{tabular}
\end{minipage}

\item \begin{minipage}[t]{\linewidth}
¿Qué afirmación es correcta?

\vspace{5pt}
\begin{tabular}{@{}p{0.47\linewidth}p{0.47\linewidth}@{}}
a)~En el metano el carbono trabaja con 3 & b)~En todos los hidruros el hidrógeno trabaja con -1 \\
c)~En un hidruro volátil el hidrógeno trabaja con +1 & d)~En los hidruros volátiles se usa la nomenclatura stock \\
\end{tabular}
\end{minipage}

\end{enumerate}
\end{document}
